\section{模型评价与推广}

\subsection{模型优点与创新}

\textbf{（1）四问共用一套可以核对的电量流骨架。}本文始终用十分钟电量连接负载、光伏、外网和电池，并统一处理90\%充电效率、90\%放电效率、储电上下限及功率限制。各问只在信息条件和费用规则发生变化时增加相应模块，避免同一物理过程被重复解释成不同含义。问题一的确定性调度由线性规划完成\cite{dantzig}，后续各问沿用同一基础关系，便于逐段复算。

\textbf{（2）把预测和高价缺口风险分开处理。}负载与光伏分别预测，使星期用电规律和近期光伏变化都能被解释\cite{hyndman}；误差准备则专门处理五倍紧急电价造成的损失不对称。候选比较显示，提升树即使改善部分预测误差，也没有超过带近期准备的完整方案。该设计说明模型追求的是最终运行费用和供电可靠性，而不是单一的拟合指标。

\textbf{（3）把日内信息的价值拆成两个可说明的部分。}问题三没有把197.28万元的总差额全部归功于新预报，而是先计算日内反馈与改单节省的175.51万元，再计算后发官方预报额外节省的21.78万元。这样可以直接回答题目关于日内预报是否值得引入的问题，也避免把权限变化误写成预测方法的贡献。

\textbf{（4）根据决策方式选择价格模型。}问题四没有机械地给两条分支使用同一种算法。零点分支采用提升树表达非线性关系\cite{friedman}；日内分支采用岭回归稳定处理高度相关的历史价格特征\cite{hoerl}。后者以极小费用差换取更少紧急购电和更清楚的系数含义，体现了经济性、风险和可解释性的综合选择。

\textbf{（5）结论能够追溯到实际运行过程。}本文使用SciPy和scikit-learn完成数值计算与候选预测\cite{virtanen,pedregosa}，线性规划由HiGHS求解\cite{huangfu}。全年费用均由保存的普通订单、紧急购电和储电轨迹重新计算，并逐段检查供需、库存、功率和结算关系。指定日期的图表进一步说明电池何时补缺、何时达到下限以及为什么仍会发生紧急购电，使模型结论既有总量证据，也有运行过程支撑。

\subsection{模型局限与改进方向}

首先，历史窗口、误差分位、日期衰减和价格模型参数主要在2025年数据内比较。现有结果能够说明本题数据范围内的取舍，但不能证明同一组参数在其他年份或其他地区仍然最优。若获得跨年数据，可按年份滚动回测，并在不同季节分别校准参数，以检验真正的跨年泛化能力。

其次，本文先预测负载、光伏和价格，再把预测结果送入优化。这样的分层结构容易解释，也便于定位误差来源，但没有同时描述三类不确定量可能共同出现的极端情景。后续可构造少量有现实含义的联合情景，比较连续阴天、高负载和高电价叠加时的计划，进一步控制尾部风险。

再次，题目没有提供电池老化费用、外送电价、网络损耗和可中断负载信息，因此本文未把这些因素写入目标函数。若用于真实微网，可增加循环寿命成本、需求响应和售电约束；末端库存的参考价值也应根据下一日市场规则重新估计，而不能直接视为可出售资产。

\subsection{模型推广性评估}

模型的核心输入只有分时负载、本地可再生能源、设备参数、交易价格和允许的决策时点。替换这些数据并重新校准预测与准备参数后，可用于社区、校园、工业园区及偏远地区的微网运行，也可扩展到多台储能、风光互补、充电站和可调负载。滚动求解方式允许系统随着新数据修正未来计划，适合部署在需要持续更新的能源管理平台中。

从社会应用看，该方法优先消纳本地光伏，并通过储能和合理购电减少高峰期对外网的额外压力，有助于提高可再生能源利用率、降低能源浪费，服务碳达峰和碳中和目标\cite{carbonplan,energyplan}。推广时还应把安全供电和居民基本用电放在成本优化之前，使节费收益建立在可靠、公平的能源服务之上。这样的应用方向既符合绿色发展和资源节约要求，也体现以人民为中心、共建共享的发展取向。

\newpage
\begin{thebibliography}{99}
\bibitem{dantzig} Dantzig G B. Linear Programming and Extensions[M]. Princeton: Princeton University Press, 1963.
\bibitem{hoerl} Hoerl A E, Kennard R W. Ridge regression: Biased estimation for nonorthogonal problems[J]. Technometrics, 1970, 12(1): 55--67.
\bibitem{friedman} Friedman J H. Greedy function approximation: A gradient boosting machine[J]. The Annals of Statistics, 2001, 29(5): 1189--1232.
\bibitem{huangfu} Huangfu Q, Hall J A J. Parallelizing the dual revised simplex method[J]. Mathematical Programming Computation, 2018, 10(1): 119--142.
\bibitem{hyndman} Hyndman R J, Athanasopoulos G. Forecasting: Principles and Practice[M]. 3rd ed. Melbourne: OTexts, 2021.
\bibitem{virtanen} Virtanen P, Gommers R, Oliphant T E, et al. SciPy 1.0: Fundamental algorithms for scientific computing in Python[J]. Nature Methods, 2020, 17: 261--272.
\bibitem{pedregosa} Pedregosa F, Varoquaux G, Gramfort A, et al. Scikit-learn: Machine learning in Python[J]. Journal of Machine Learning Research, 2011, 12: 2825--2830.
\bibitem{carbonplan} 国务院. 2030年前碳达峰行动方案[Z]. 2021.
\bibitem{energyplan} 国家发展改革委, 国家能源局. “十四五”现代能源体系规划[Z]. 2022.
\end{thebibliography}
